\numoptstrue
\yearpaper{2025}{Day 1 \textperiodcentered{} First Session \textperiodcentered{} Version A1 \textperiodcentered{} 60 Questions \textperiodcentered{} 60 Marks \textperiodcentered{} 80 Minutes}{}
\begin{multicols}{2}

\QQ{Which of the following expressions can be deduced on the basis of dimensional analysis? (All symbols have their usual meanings)\ntag{A}{9}{keph101}{9}{}{Class XI Part 1 \textperiodcentered{} Ch 1}\nsrc{None found in NEET/AIPMT/AIEEE/JEE papers, 2005--2026 (NCERT concept application)}}
{\opt{$s = ut+\frac12at^2$}{$x = A\cos\omega t$}{$N = N_0e^{-\lambda t}$}{$F = 6\pi\eta rv$}}

\QQ{Two stones begin to fall from rest from the same height, with the second stone starting to fall `$\Delta t$' seconds after the first falls from rest. The distance of separation between the two stones becomes `$H$', `$t_0$' seconds after the first stone starts its motion. Then $t_0$ is equal to\ntag{A}{18}{keph102}{6}{}{Class XI Part 1 \textperiodcentered{} Ch 2}\nsrc{None found in NEET/AIPMT/AIEEE/JEE papers, 2005--2026 (NCERT concept application)}}
{\opt{$\dfrac{H}{g\Delta t}$}{$\dfrac{H}{\Delta t}+\dfrac{\Delta t}{2g}$}{$\dfrac{H}{g\Delta t}-\dfrac{\Delta t}{2}$}{$\dfrac{H}{g\Delta t}+\dfrac{\Delta t}{2}$}}

\QQ{In the projectile motion of a particle on a level ground, which of the following remains constant with reference to time and position?\ntag{D}{38}{keph103}{12}{}{Class XI Part 1 \textperiodcentered{} Ch 3}}
{\opt{Vertical component of the velocity of the projectile}{Average velocity between any two points on the path}{Horizontal component of velocity}{Angle between the instantaneous velocity with the horizontal}}

\QQ{A particle is in uniform circular motion. The equation of its trajectory is given by $(x-2)^2+y^2 = 25$, where $x$ and $y$ are in meter. The speed of the particle is $2\unit{m\,s^{-1}}$. When the particle attains the lowest `$y$' co-ordinate, the acceleration of the particle is (in $\mathrm{m\,s^{-2}}$)\ntag{D}{41}{keph103}{15}{}{Class XI Part 1 \textperiodcentered{} Ch 3}}
{\opt{$0.4\,\hat i$}{$0.4\,\hat j$}{$0.8\,\hat i$}{$0.8\,\hat j$}}

\QQ{A wooden block of mass M lies on a rough floor. Another wooden block of the same mass is hanging from the point O through strings as shown in the figure. To achieve equilibrium, the co-efficient of static friction between the block on the floor with the floor itself is
\fig{\begin{tikzpicture}[scale=0.8]
\fill[pattern=north east lines] (-0.3,-0.12) rectangle (1.3,0);
\draw (-0.3,0) -- (1.3,0);
\draw (0,0) rectangle (0.8,0.6) node[midway,lbl]{M};
\draw (0.8,0.3) -- (2.6,0.3);
\fill (2.6,0.3) circle (1.2pt) node[above,lbl]{O};
\draw (2.6,0.3) -- (4.0,1.1);
\fill[pattern=north east lines] (4.0,0.5) rectangle (4.12,1.9); \draw (4.0,0.5)--(4.0,1.9);
\draw[dashed] (2.6,0.3) -- (3.5,0.3);
\draw (3.1,0.3) arc (0:29:0.5); \node[lbl] at (3.3,0.47) {$\theta$};
\draw (2.6,0.3) -- (2.6,-0.6);
\draw (2.3,-0.6) rectangle (2.9,-1.15) node[midway,lbl]{M};
\end{tikzpicture}}\ntag{A}{61}{keph104}{13}{}{Class XI Part 1 \textperiodcentered{} Ch 4}\nsrc{None found in NEET/AIPMT/AIEEE/JEE papers, 2005--2026 (NCERT concept application)}}
{\opt{$\mu = \cos\theta$}{$\mu = \cot\theta$}{$\mu = \sin\theta$}{$\mu = \tan\theta$}}

\QQ{A block of certain mass is placed on a rough floor. The coefficients of static and kinetic friction between the block and the floor are 0.4 and 0.25 respectively. A constant horizontal force $F = 20\unit{N}$ acts on it so that the velocity of the block varies with time according to the following graph. The mass of the block is nearly (Take $g\simeq10\unit{m\,s^{-2}}$)
\fig{\begin{tikzpicture}[xscale=0.55,yscale=0.09]
\foreach \x in {1,...,5} \draw[densely dotted] (\x,0) -- (\x,22);
\foreach \y in {5,10,15,20} {\draw[densely dotted] (0,\y) -- (5.5,\y); \node[lbl,left] at (0,\y) {\y};}
\draw[->] (0,0) -- (5.9,0) node[right,lbl]{$t$ (s)};
\draw[->] (0,0) -- (0,25) node[above,lbl]{$v$ (m/s)};
\foreach \x in {0,...,5} \node[lbl,below] at (\x,0) {\x};
\draw[line width=0.9pt] (0,0) -- (3,20);
\end{tikzpicture}}\ntag{A}{61}{keph104}{13}{}{Class XI Part 1 \textperiodcentered{} Ch 4}\nsrc{None found in NEET/AIPMT/AIEEE/JEE papers, 2005--2026 (NCERT concept application)}}
{\opt{$2.2\unit{kg}$}{$4.4\unit{kg}$}{$1.2\unit{kg}$}{$1.0\unit{kg}$}}

\QQ{A body of mass $0.25\unit{kg}$ travels along a straight line from $x=0$ to $x=2\unit{m}$ with a speed $v = kx^{3/2}$ where $k = 2$ SI units. The work done by the net force during this displacement is\ntag{D}{73}{keph105}{3}{}{Class XI Part 1 \textperiodcentered{} Ch 5}}
{\opt{$4\unit{J}$}{$8\unit{J}$}{$16\unit{J}$}{$32\unit{J}$}}

\QQ{During an elastic collision between two bodies, which of the following statements are correct?
\begin{stmts}
\item The initial kinetic energy is equal to the final kinetic energy of the system.
\item The linear momentum is conserved.
\item The kinetic energy during $\Delta t$ (the collision time) is not conserved.
\end{stmts}\ntag{D}{84}{keph105}{14}{}{Class XI Part 1 \textperiodcentered{} Ch 5}}
{\opt{I and II only}{II and III only}{I and III only}{I, II and III}}

\QQ{Three particles of mass $1\unit{kg}$, $2\unit{kg}$ and $3\unit{kg}$ are placed at the vertices A, B and C respectively of an equilateral triangle ABC of side $1\unit{m}$. The centre of mass of the system from vertex A (located at origin) is\ntag{D}{97}{keph106}{6}{ \textperiodcentered{} Example 6.1}{Class XI Part 1 \textperiodcentered{} Ch 6}}
{\opt{$(0,0)$}{$\left(\dfrac{7}{12},\dfrac{3\sqrt3}{12}\right)$}{$\left(\dfrac{9}{12},\dfrac{3\sqrt3}{12}\right)$}{$\left(\dfrac{7}{12},\dfrac{6+3\sqrt3}{12}\right)$}}

\QQ{Two fly wheels are connected by a non-slipping belt as shown in the figure. $I_1 = 4\unit{kg\,m^2}$, $r_1 = 20\unit{cm}$, $I_2 = 20\unit{kg\,m^2}$ and $r_2 = 30\unit{cm}$. A torque of $10\unit{N\,m}$ is applied on the smaller wheel. Then match the entries of column I with appropriate entries of column II.\par\smallskip
{\small\begin{tabular}{@{}|l|p{0.44\linewidth}|l|c|@{}}\hline
I & Quantities & II & Value (SI)\\\hline
(a) & Angular acceleration of smaller wheel & (i) & $\frac53$\\[3pt]\hline
(b) & Torque on the larger wheel & (ii) & $\frac{100}{3}$\\[3pt]\hline
(c) & Angular acceleration of larger wheel & (iii) & $\frac52$\\[3pt]\hline
\end{tabular}}
\fig{\begin{tikzpicture}[scale=0.75]
\draw (0,0) circle (0.8); \draw (2.3,0) circle (1.2);
\draw ($(0,0)+(105:0.8)$) -- ($(2.3,0)+(97:1.2)$);
\draw ($(0,0)+(-105:0.8)$) -- ($(2.3,0)+(-97:1.2)$);
\draw[->] (0,0) -- (55:0.8) node[pos=0.6,right,lbl]{$r_1$};
\draw[->] (2.3,0) -- ($(2.3,0)+(120:1.2)$) node[pos=0.6,right,lbl]{$r_2$};
\fill (0,0) circle (1.5pt) (2.3,0) circle (1.5pt);
\end{tikzpicture}}\ntag{A}{119}{keph106}{28}{}{Class XI Part 1 \textperiodcentered{} Ch 6}\nsrc{None found in NEET/AIPMT/AIEEE/JEE papers, 2005--2026 (NCERT concept application)}}
{\opt{a--iii, b--ii, c--i}{a--ii, b--iii, c--i}{a--iii, b--i, c--ii}{a--ii, b--i, c--iii}}

\QQ{If $r_p, v_p, L_p$ and $r_a, v_a, L_a$ are radii, velocities and angular momenta of a planet at perihelion and aphelion of its elliptical orbit around the Sun respectively, then\ntag{D}{129}{keph107}{3}{ \textperiodcentered{} Example 7.1}{Class XI Part 1 \textperiodcentered{} Ch 7}}
{\opt{$r_p<r_a,\ v_p<v_a,\ L_a<L_p$}{$r_p>r_a,\ v_p>v_a,\ L_a>L_p$}{$r_p<r_a,\ v_p>v_a,\ L_a=L_p$}{$r_p>r_a,\ v_p<v_a,\ L_a=L_p$}}

\QQ{The total energy of a satellite in a circular orbit at a distance $(R+h)$ from the center of the Earth varies as [$R$ is the radius of the Earth and $h$ is the height of the orbit from Earth's surface]\ntag{D}{138}{keph107}{12}{}{Class XI Part 1 \textperiodcentered{} Ch 7}}
{\opt{$\dfrac{1}{R+h}$}{$-\dfrac{1}{R+h}$}{$\dfrac{1}{(R+h)^2}$}{$-\dfrac{1}{(R+h)^2}$}}

\QQ{Two wires A and B are made of same material. Their diameters are in the ratio of 1:2 and lengths are in the ratio of 1:3. If they are stretched by the same force, then increase in their lengths will be in the ratio of\ntag{D}{170}{keph201}{4}{}{Class XI Part 2 \textperiodcentered{} Ch 8}}
{\opt{4:3}{3:4}{2:3}{3:2}}

\QQ{A horizontal pipe carries water in a streamlined flow. At a point along the pipe, where the cross-sectional area is $10\unit{cm^2}$, the velocity of water is $1\unit{m\,s^{-1}}$ and the pressure is $2000\unit{Pa}$. What is the pressure of water at another point where the cross-sectional area is $5\unit{cm^2}$? [Density of water $= 1000\unit{kg\,m^{-3}}$]\ntag{A}{188}{keph202}{9}{}{Class XI Part 2 \textperiodcentered{} Ch 9}\nsrc{JEE Main 2025 (23 Jan, Shift 2) -- water in a horizontal pipe, pressure readings (Bernoulli)}}
{\opt{$200\unit{Pa}$}{$300\unit{Pa}$}{$400\unit{Pa}$}{$500\unit{Pa}$}}

\QQ{Three metal rods of the same material and identical in all respects are joined as shown in the figure. The temperatures at the ends of these rods are maintained as indicated. Assuming no heat energy loss occurs through the curved surfaces of the rods, the temperature at the junction $x$ is
\fig{\begin{tikzpicture}[scale=0.8]
\draw[line width=2.2pt] (0,0) -- (1.8,0); \draw[line width=2.2pt] (1.8,0) -- (3.2,1); \draw[line width=2.2pt] (1.8,0) -- (3.2,-1);
\draw[parchment,line width=1.2pt] (0,0) -- (1.8,0) -- (3.2,1) (1.8,0) -- (3.2,-1);
\node[lbl,left] at (0,0) {$0^\circ$C}; \node[lbl,right] at (3.2,1) {$90^\circ$C}; \node[lbl,right] at (3.2,-1) {$90^\circ$C};
\node[lbl,above] at (1.8,0.05) {$x$};
\end{tikzpicture}}\ntag{A}{214}{keph203}{13}{}{Class XI Part 2 \textperiodcentered{} Ch 10}\nsrc{NEET 2025 -- three heat-conducting rods joined, temperature at the junction}}
{\opt{$45^\circ$C}{$60^\circ$C}{$30^\circ$C}{$20^\circ$C}}

\QQ{A gas is taken from state A to state B along two different paths 1 and 2. The heat absorbed and work done by the system along these two paths are $Q_1$ and $Q_2$ and $W_1$ and $W_2$ respectively. Then\ntag{D}{229}{keph204}{4}{}{Class XI Part 2 \textperiodcentered{} Ch 11}}
{\opt{$Q_1 = Q_2$}{$W_1 = W_2$}{$Q_1-W_1 = Q_2-W_2$}{$Q_1+W_1 = Q_2+W_2$}}

\QQ{At $27^\circ$C temperature, the mean kinetic energy of the atoms of an ideal gas is $E_1$. If the temperature is increased to $327^\circ$C, then the mean kinetic energy of the atoms will be\ntag{D}{250}{keph205}{7}{}{Class XI Part 2 \textperiodcentered{} Ch 12}}
{\opt{$\dfrac{E_1}{2}$}{$\dfrac{E_1}{\sqrt2}$}{$\sqrt2\,E_1$}{$2E_1$}}

\QQ{The variations of kinetic energy $K(x)$, potential energy $U(x)$ and total energy as a function of displacement of a particle in SHM is as shown in the figure. The value of $|x_0|$ is
\fig{\begin{tikzpicture}[scale=0.85]
\draw[->,densely dashed] (-2.4,0) -- (2.6,0) node[right,lbl]{$x$};
\draw[->,densely dashed] (0,0) -- (0,2.2) node[above,lbl]{$E,U,K$};
\draw (-2,1.6) -- (2,1.6); \node[lbl,above right] at (2,1.6) {$E$};
\draw[domain=-2:2,samples=40] plot (\x,{0.4*\x*\x}); \node[lbl,right] at (1.75,1.0) {$U(x)$};
\draw[domain=-2:2,samples=40] plot (\x,{1.6-0.4*\x*\x}); \node[lbl,right] at (1.9,0.35) {$K(x)$};
\draw[densely dotted] (-2,0) -- (-2,1.6) (2,0) -- (2,1.6) (-1.414,0) -- (-1.414,0.8) (1.414,0) -- (1.414,0.8);
\node[lbl,below] at (-2.05,0) {$-A$}; \node[lbl,below] at (-1.35,0) {$-x_0$};
\node[lbl,below] at (1.35,0) {$x_0$}; \node[lbl,below] at (2.05,0) {$+A$};
\end{tikzpicture}}\ntag{A}{268}{keph206}{10}{}{Class XI Part 2 \textperiodcentered{} Ch 13}\nsrc{JEE Main 2021 (17 Mar, Shift 1) -- displacement where KE and PE of SHM are equal}}
{\opt{$\dfrac{A}{2}$}{$2A$}{$\dfrac{A}{\sqrt2}$}{$\sqrt2\,A$}}

\QQ{The angle between the particle velocity and wave velocity in a transverse wave is [except when the particle passes through the mean position]\ntag{D}{280}{keph207}{3}{}{Class XI Part 2 \textperiodcentered{} Ch 14}}
{\opt{Zero radian}{$\dfrac{\pi}{4}$ radian}{$\dfrac{\pi}{2}$ radian}{$\pi$ radian}}

\QQ{A metallic sphere of radius $R$ carrying a charge $q$ is kept at certain distance from another metallic sphere of radius $R/4$ carrying a charge $Q$. What is the electric flux at any point inside the metallic sphere of radius $R$ due to the sphere of radius $R/4$?
\fig{\begin{tikzpicture}[scale=0.7]
\draw (0,0) circle (0.8); \draw[->] (0,0) -- (30:0.8) node[pos=0.5,above,lbl]{$R$}; \fill (0,0) circle (1.2pt);
\foreach \a in {0,30,...,330} \node[font=\tiny] at (\a:0.98) {$\times$};
\node[lbl] at (0,1.25) {$q$};
\draw (2.6,0) circle (0.45); \draw[->] (2.6,0) -- ($(2.6,0)+(-20:0.45)$); \fill (2.6,0) circle (1.2pt);
\node[lbl] at (3.45,-0.25) {$\frac R4$};
\foreach \a in {0,36,...,324} \node[font=\tiny] at ($(2.6,0)+(\a:0.62)$) {$\times$};
\node[lbl] at (2.6,0.9) {$Q$};
\end{tikzpicture}}\ntag{A}{62}{leph102}{18}{}{Class XII Part 1 \textperiodcentered{} Ch 2}\nsrc{None found in NEET/AIPMT/AIEEE/JEE papers, 2005--2026 (NCERT concept application)}}
{\opt{$\dfrac{Q}{\varepsilon_0}$}{$\dfrac{Q}{\varepsilon_0}-\dfrac{q}{\varepsilon_0}$}{Zero}{$\dfrac{q}{\varepsilon_0}-\dfrac{Q}{\varepsilon_0}$}}

\QQ{You are given a dipole of charge $+q$ and $-q$ separated by a distance $2R$. A sphere `A' of radius `$R$' passes through the centre of the dipole as shown below and another sphere `B' of radius `$2R$' passes through the charge $+q$. Then the electric flux through the sphere A is
\fig{\begin{tikzpicture}[scale=0.55]
\draw (0,0) circle (1); \draw (0,0) circle (2);
\fill (0,0) circle (2pt) node[above left,lbl]{$-q$};
\fill (2,0) circle (2pt) node[above left,lbl]{$+q$};
\draw[densely dashed] (0,0) -- (2,0);
\draw[<->] (0,-0.3) -- (2,-0.3) node[midway,below,lbl]{$2R$};
\node[lbl] at (0.75,1.0) {A}; \node[lbl] at (1.6,1.6) {B};
\end{tikzpicture}}\ntag{A}{29}{leph101}{29}{}{Class XII Part 1 \textperiodcentered{} Ch 1}\nsrc{NEET 2019 -- sphere enclosing an electric dipole, total flux; also JEE Main 2021 (26 Feb, Shift 2)}}
{\opt{$-q/\varepsilon_0$}{$q/\varepsilon_0$}{Zero}{$2q/\varepsilon_0$}}

\QQ{A potential at a point A is $-3\unit{V}$ and that at another point B is $5\unit{V}$. What is the work done in carrying a charge of $5\unit{mC}$ from B to A?\ntag{D}{46}{leph102}{2}{}{Class XII Part 1 \textperiodcentered{} Ch 2}}
{\opt{$-40\unit{J}$}{$-0.04\unit{J}$}{$-0.4\unit{J}$}{$-4\unit{J}$}}

\QQ{Charges are uniformly spread on the surface of a conducting sphere. The electric field from the centre of sphere to a point outside the sphere varies with distance $r$ from the centre as\ntag{D}{35}{leph101}{35}{}{Class XII Part 1 \textperiodcentered{} Ch 1}}
{\optfig{\gaxes{$r$}{$E$}{\draw (0,0) -- (0.9,1.6) .. controls (1.0,0.6) and (1.3,0.4) .. (2,0.35); \draw[densely dashed] (0.9,1.6) -- (0.9,0);}}
{\gaxes{$r$}{$E$}{\draw (0,0) -- (0.9,0); \draw (0.9,1.6) .. controls (1.0,0.6) and (1.3,0.4) .. (2,0.35); \draw[densely dashed] (0.9,1.6) -- (0.9,0);}}
{\gaxes{$r$}{$E$}{\draw (0,1.5) -- (1.9,0);}}
{\gaxes{$r$}{$E$}{\draw (0.1,1.6) .. controls (0.5,0.4) and (1.1,0.3) .. (2,0.3);}}}

\QQ{Match Column-I with Column-II related to an electric dipole of dipole moment $\vec p$ that is placed in a uniform electric field $\vec E$.\par\smallskip
\begin{adjustbox}{max width=\linewidth}\begin{tabular}{@{}ll@{\quad}ll@{}}
\multicolumn{2}{@{}l}{\textit{Column-I} (angle between $\vec p$ and $\vec E$)} & \multicolumn{2}{l}{\textit{Column-II} (P.E.)}\\[2pt]
(a) & $180^\circ$ & (i) & $-pE$\\
(b) & $120^\circ$ & (ii) & $pE$\\
(c) & $90^\circ$ & (iii) & $\frac12pE$\\
 & & (iv) & Zero
\end{tabular}\end{adjustbox}\ntag{D}{60}{leph102}{16}{}{Class XII Part 1 \textperiodcentered{} Ch 2}}
{\opt{a--ii, b--iii, c--iv}{a--i, b--ii, c--iii}{a--ii, b--iii, c--i}{a--ii, b--i, c--iv}}

\QQ{Which of the following statements is not true?\ntag{D}{54}{leph102}{10}{}{Class XII Part 1 \textperiodcentered{} Ch 2}}
{\opt{Electric field is always perpendicular to an equipotential surface.}{Work done to move a charge on an equipotential surface is not zero}{Equipotential surfaces are the surfaces where the potential is constant}{Equipotential surfaces for a uniform electric field are parallel and equidistant from each other.}}

\QQ{Which of the following is a correct statement?\ntag{D}{30}{leph101}{30}{}{Class XII Part 1 \textperiodcentered{} Ch 1}}
{\opt{Gauss's law is true for any closed surface}{Gauss's law is true for any open surface}{Gauss's law is not applicable when charges are not symmetrically distributed over a closed surface.}{Gauss's law does not hold good for a charge situated outside the Gaussian surface.}}

\QQ{In the following circuit, the terminal voltage across the cell is
\fig{\begin{circuitikz}[scale=0.7,transform shape]
\draw (0,1.6) to[battery1,l=$2\,\mathrm{V}$] (2,1.6) to[R,l=$0.1\,\Omega$] (4,1.6) -- (4,0) to[R,l=$3.9\,\Omega$] (0,0) -- (0,1.6);
\end{circuitikz}}\ntag{A}{95}{leph103}{15}{}{Class XII Part 1 \textperiodcentered{} Ch 3}\nsrc{NEET 2024 -- terminal voltage of a battery with internal resistance}}
{\opt{$0.52\unit{V}$}{$1.68\unit{V}$}{$1.95\unit{V}$}{$2.71\unit{V}$}}

\QQ{Two cells of emfs $E_1$ and $E_2$ and internal resistances $r_1$ and $r_2$ ($E_2>E_1$ and $r_2>r_1$) respectively, are connected in parallel as shown in figure. The equivalent emf of the combination is $E_{eq}$. Then
\fig{\begin{circuitikz}[scale=0.7,transform shape]
\draw (-0.8,1) node[circ]{} node[left]{A} -- (0,1);
\draw (0,1) -- (0,2) to[battery1,l=$E_1$] (1.8,2) to[R,l=$r_1$] (3.6,2) -- (3.6,1);
\draw (0,1) -- (0,0) to[battery1,l_=$E_2$] (1.8,0) to[R,l_=$r_2$] (3.6,0) -- (3.6,1);
\draw (3.6,1) -- (4.4,1) node[circ]{} node[right]{B};
\end{circuitikz}}\ntag{A}{96}{leph103}{16}{}{Class XII Part 1 \textperiodcentered{} Ch 3}\nsrc{None found in NEET/AIPMT/AIEEE/JEE papers, 2005--2026 (NCERT concept application)}}
{\opt{$E_1<E_{eq}<E_2$ and $E_{eq}$ is nearer $E_1$}{$E_1<E_{eq}<E_2$ and $E_{eq}$ is nearer $E_2$}{$E_{eq}>E_2$}{$E_{eq}<E_1$}}

\QQ{The variations of resistivity $\rho$ with absolute temperature $T$ for three different materials X, Y and Z are shown in the graph below. Identify the materials X, Y and Z.
\fig{\begin{tikzpicture}[scale=0.8]
\draw[->] (-0.2,0) -- (3.3,0) node[right,lbl]{$T$};
\draw[->] (0,-0.2) -- (0,2.8) node[above,lbl]{$\rho$};
\draw (0,1.0) -- (2.9,1.8) node[right,lbl]{Y};
\draw (0,0.55) .. controls (1.5,0.55) and (2.0,1.2) .. (2.25,2.5) node[above,lbl]{X};
\draw (0.3,2.5) node[above,lbl]{Z} .. controls (0.8,0.5) and (1.5,0.2) .. (2.7,0.2);
\end{tikzpicture}}\ntag{D}{90}{leph103}{10}{}{Class XII Part 1 \textperiodcentered{} Ch 3}}
{\opt{X -- copper, Y -- nichrome, Z -- semiconductor}{X -- copper, Y -- semiconductor, Z -- nichrome}{X -- semiconductor, Y -- nichrome, Z -- copper}{X -- nichrome, Y -- copper, Z -- semiconductor}}

\QQ{Given, a current carrying wire of non-uniform cross-section, which of the following is constant throughout the length of wire?\ntag{A}{82}{leph103}{2}{}{Class XII Part 1 \textperiodcentered{} Ch 3}\nsrc{AIPMT 2015 -- conductor of non-uniform cross-section, quantity that remains constant}}
{\opt{Current, electric field and drift speed}{Drift speed}{Current and drift speed}{Current only}}

\QQ{The graph between variation of resistance of a metal wire as a function of its diameter keeping other parameters like length and temperature constant is\ntag{D}{84}{leph103}{4}{}{Class XII Part 1 \textperiodcentered{} Ch 3}}
{\optfig{\gaxes{$D$}{$R$}{\draw (0,0) -- (1.5,1.5);}}
{\gaxes{$D$}{$R$}{\draw (0,1.5) -- (1.4,0.1);}}
{\gaxes{$D$}{$R$}{\draw[domain=0:1.7] plot (\x,{0.5*\x*\x});}}
{\gaxes{$D$}{$R$}{\draw[domain=0.3:1.9,samples=40] plot (\x,{0.14/(\x*\x)});}}}

\QQ{Two thin long parallel wires separated by a distance `$r$' from each other in vacuum carry a current of $I$ ampere in opposite directions. Then, they will\ntag{D}{123}{leph104}{17}{}{Class XII Part 1 \textperiodcentered{} Ch 4}}
{\opt{Attract each other with a force per unit length of $\dfrac{\mu_0I^2}{2\pi r^2}$}{Attract each other with a force per unit length of $\dfrac{\mu_0I^2}{2\pi r}$}{Repel each other with a force per unit length of $\dfrac{\mu_0I^2}{2\pi r}$}{Repel each other with a force per unit length of $\dfrac{\mu_0I^2}{2\pi r^2}$}}

\QQ{A solenoid is $1\unit{m}$ long and $4\unit{cm}$ in diameter. It has five layers of windings of 1000 turns each and carries a current of $7\unit{A}$. The magnetic field at the centre of the solenoid is\ntag{D}{135}{leph104}{29}{ \textperiodcentered{} Exercise 4.8}{Class XII Part 1 \textperiodcentered{} Ch 4}}
{\opt{$439.6\unit{T}$}{$0.4396\times10^{-5}\unit{T}$}{$4.396\times10^{-2}\unit{T}$}{$43.96\times10^{-2}\unit{T}$}}

\QQ{Two similar galvanometers are converted into an ammeter and a milliammeter. The shunt resistance of ammeter as compared to the shunt resistance of milliammeter will be\ntag{A}{130}{leph104}{24}{}{Class XII Part 1 \textperiodcentered{} Ch 4}\nsrc{None found in NEET/AIPMT/AIEEE/JEE papers, 2005--2026 (NCERT concept application)}}
{\opt{Equal}{Zero}{More}{Less}}

\QQ{Which of the following statements is true in respect of diamagnetic substances?\ntag{D}{147}{leph105}{12}{}{Class XII Part 1 \textperiodcentered{} Ch 5}}
{\opt{Susceptibility is small and negative}{They are feebly attracted by magnets}{Permeability is greater than 1000}{Susceptibility decreases with temperature.}}

\QQ{Identify the correct statement\ntag{D}{119}{leph104}{13}{}{Class XII Part 1 \textperiodcentered{} Ch 4}}
{\opt{The magnetic field inside a solenoid is non-uniform}{A current carrying conductor produces an electric field around it.}{A straight current carrying conductor has circular magnetic field lines around it.}{The direction of magnetic field due to a current element is given by Fleming's Left Hand Rule}}

\QQ{Which of the following graphs represents the variation of magnetic field $B$ with perpendicular distance `$r$' from an infinitely long, straight conductor carrying current?\ntag{D}{118}{leph104}{12}{}{Class XII Part 1 \textperiodcentered{} Ch 4}}
{\optfig{\gaxes{$r$}{$B$}{\draw (0,1.6) -- (1.5,0);}}
{\gaxes{$r$}{$B$}{\draw (0,0) -- (1.6,1.5);}}
{\gaxes{$r$}{$B$}{\draw[domain=0.18:1.9,samples=40] plot (\x,{0.28/\x});}}
{\gaxes{$r$}{$B$}{\draw (0,0) .. controls (0.4,1.2) and (1.2,1.7) .. (1.9,0.9);}}}

\QQ{If we consider an electron and a photon of same de-Broglie wavelength, then they will have same\ntag{D}{285}{leph203}{12}{}{Class XII Part 2 \textperiodcentered{} Ch 11}}
{\opt{Momentum}{Angular momentum}{Energy}{Velocity}}

\QQ{The anode voltage of a photocell is kept fixed. The frequency of the light falling on the cathode is gradually increased. Then the correct graph which shows the variation of photo current $I$ with the frequency $\nu$ of incident light is\ntag{D}{279}{leph203}{6}{}{Class XII Part 2 \textperiodcentered{} Ch 11}}
{\optfig{\gaxes{$\nu$}{$I$}{\draw (0.6,0) -- (1.6,1.7); \node[lbl,below] at (0.6,0) {$\nu_0$};}}
{\gaxes{$\nu$}{$I$}{\draw[densely dashed] (0.5,0) -- (0.5,1.2); \draw (0.5,1.2) -- (1.9,1.2); \node[lbl,below] at (0.5,0) {$\nu_0$};}}
{\gaxes{$\nu$}{$I$}{\draw (0.15,0) .. controls (0.6,0.2) and (0.6,1.3) .. (1.1,1.35) -- (1.9,1.35);}}
{\gaxes{$\nu$}{$I$}{\draw (0.1,1.6) .. controls (0.5,0.6) and (1.0,0.3) .. (1.7,0.1);}}}

\QQ{When a bar magnet is pushed towards the coil, along its axis, as shown in the figure, the galvanometer pointer deflects towards X. When this magnet is pulled away from the coil, the galvanometer pointer
\fig{\begin{tikzpicture}[scale=0.7]
\draw (0,0) rectangle (2.6,1.1);
\draw (1.0,1.1) .. controls (0.6,2.0) and (0.4,3.0) .. (1.2,3.1) .. controls (2.0,3.0) and (1.8,2.0) .. (1.4,1.1);
\draw (2.9,1.9) rectangle (4.6,2.5); \node[lbl] at (3.2,2.2) {N}; \node[lbl] at (4.3,2.2) {S};
\draw[fill=parchment] (1.3,0) circle (0.22); \draw[->] (1.3,-0.12) -- (1.3,0.15); \node[lbl,below] at (1.3,-0.25) {G};
\node[lbl] at (-0.3,-0.25) {X$'$}; \node[lbl] at (2.9,-0.25) {X};
\end{tikzpicture}}\ntag{D}{155}{leph106}{2}{}{Class XII Part 1 \textperiodcentered{} Ch 6}}
{\opt{deflects towards X}{deflects towards X$'$}{does not deflect}{oscillates}}

\QQ{A square loop of side $2\unit{m}$ lies in the Y-Z plane in a region having a magnetic field $\vec B = (5\hat i+3\hat j-4\hat k)\unit{T}$. The magnitude of magnetic flux through the square loop is\ntag{D}{157}{leph106}{4}{}{Class XII Part 1 \textperiodcentered{} Ch 6}}
{\opt{$10\unit{Wb}$}{$20\unit{Wb}$}{$12\unit{Wb}$}{$16\unit{Wb}$}}

\QQ{In domestic electric mains supply, the voltage and the current are\ntag{D}{177}{leph107}{1}{}{Class XII Part 1 \textperiodcentered{} Ch 7}}
{\opt{AC voltage and AC current}{AC voltage and DC current}{DC voltage and DC current}{DC voltage and AC current}}

\QQ{A sinusoidal voltage produced by an AC generator at any instant $t$ is given by an equation $V = 311\sin314t$. The rms value of voltage and frequency are respectively\ntag{D}{180}{leph107}{4}{}{Class XII Part 1 \textperiodcentered{} Ch 7}}
{\opt{$200\unit{V}$, $100\unit{Hz}$}{$200\unit{V}$, $50\unit{Hz}$}{$220\unit{V}$, $100\unit{Hz}$}{$220\unit{V}$, $50\unit{Hz}$}}

\QQ{A series LCR circuit containing an AC source of $100\unit{V}$ has an inductor and a capacitor of reactances $24\unit{\Omega}$ and $16\unit{\Omega}$ respectively. If a resistance of $6\unit{\Omega}$ is connected in series, then the potential difference across the series combination of inductor and capacitor only is\ntag{A}{186}{leph107}{10}{}{Class XII Part 1 \textperiodcentered{} Ch 7}\nsrc{None found in NEET/AIPMT/AIEEE/JEE papers, 2005--2026 (NCERT concept application)}}
{\opt{$40\unit{V}$}{$80\unit{V}$}{$400\unit{V}$}{$8\unit{V}$}}

\QQ{Match the following types of waves with their wavelength ranges\par\smallskip
\begin{tabular}{@{}ll@{\quad}ll@{}}
\multicolumn{2}{@{}l}{\textit{Waves}} & \multicolumn{2}{l}{\textit{Wavelength ranges}}\\[2pt]
i. & Microwave & (a) & $700\unit{nm}$ to $400\unit{nm}$\\
ii. & Visible light & (b) & $1\unit{nm}$ to $10^{-3}\unit{nm}$\\
iii. & Ultraviolet & (c) & $0.1\unit{m}$ to $1\unit{mm}$\\
iv. & X-rays & (d) & $400\unit{nm}$ to $1\unit{nm}$
\end{tabular}\ntag{D}{211}{leph108}{11}{}{Class XII Part 1 \textperiodcentered{} Ch 8}}
{\opt{i--a, ii--d, iii--b, iv--c}{i--c, ii--a, iii--d, iv--b}{i--d, ii--b, iii--c, iv--a}{i--b, ii--c, iii--a, iv--d}}

\QQ{A ray of light passes from vacuum into a medium of refractive index $n$. If the angle of incidence is twice the angle of refraction, then the angle of incidence in terms of refractive index is\ntag{A}{228}{leph201}{8}{}{Class XII Part 2 \textperiodcentered{} Ch 9}\nsrc{None found in NEET/AIPMT/AIEEE/JEE papers, 2005--2026 (NCERT concept application)}}
{\opt{$\cos^{-1}\!\left(\dfrac n2\right)$}{$\sin^{-1}\!\left(\dfrac n2\right)$}{$2\cos^{-1}\!\left(\dfrac n2\right)$}{$2\sin^{-1}\!\left(\dfrac n2\right)$}}

\QQ{A convex lens has power $P$. It is cut into two halves along its principal axis. Further one piece (out of two halves) is cut into two halves perpendicular to the principal axis as shown in figure. Choose the incorrect option for the reported lens pieces.
\fig{\begin{tikzpicture}[scale=0.75]
\draw (0,1.2) .. controls (0.45,0.4) and (0.45,-0.4) .. (0,-1.2) .. controls (-0.45,-0.4) and (-0.45,0.4) .. (0,1.2);
\draw[dashed] (-1.2,0) -- (2.2,0);
\begin{scope}[xshift=3cm]
\draw (-0.5,-0.08) .. controls (-0.5,-0.7) and (-0.25,-1.2) .. (0,-1.25) .. controls (0.25,-1.2) and (0.5,-0.7) .. (0.5,-0.08) -- cycle;
\draw (-0.5,0.08) .. controls (-0.5,0.7) and (-0.25,1.2) .. (-0.05,1.25) -- (-0.05,0.08) -- cycle;
\draw (0.05,0.08) -- (0.05,1.25) .. controls (0.25,1.2) and (0.5,0.7) .. (0.5,0.08) -- cycle;
\node[lbl] at (-0.35,1.45) {$L_3$}; \node[lbl] at (0.4,1.45) {$L_2$}; \node[lbl] at (0.75,-1.1) {$L_1$};
\end{scope}
\end{tikzpicture}}\ntag{A}{236}{leph201}{16}{}{Class XII Part 2 \textperiodcentered{} Ch 9}\nsrc{JEE Main 2022 (27 Jun, Shift 2) -- same question: lens of power $P$ cut along and across the principal axis}}
{\opt{Power of $L_1$ is $\dfrac P2$}{Power of $L_2$ is $\dfrac P2$}{Power of $L_3$ is $\dfrac P2$}{Power of $L_1$ is $P$}}

\QQ{The image formed by an objective lens of a compound microscope is\ntag{D}{243}{leph201}{23}{}{Class XII Part 2 \textperiodcentered{} Ch 9}}
{\opt{Virtual and diminished}{Real and diminished}{Real and enlarged}{Virtual and enlarged}}

\QQ{If $r$ and $r'$ denote the angles inside the prism having angle of prism $50^\circ$ considering that during interval of time from $t=0$ to $t=t$, $r$ varies with time as $r = 10^\circ+t^2$. During this time $r'$ will vary with time as
\fig{\begin{tikzpicture}[scale=0.8]
\coordinate (T) at (0,2.2); \coordinate (L) at (-1.2,0); \coordinate (R) at (1.2,0);
\draw (L) -- (T) -- (R) -- cycle;
\draw (T) ++(-118:0.45) arc (-118:-62:0.45); \node[lbl] at (0,1.55) {$50^\circ$};
\coordinate (P1) at ($(L)!0.55!(T)$); \coordinate (P2) at ($(R)!0.6!(T)$);
\draw[->] ($(P1)+(-1.4,-0.55)$) -- ($(P1)+(-0.6,-0.235)$); \draw ($(P1)+(-0.6,-0.235)$) -- (P1);
\draw (P1) -- (P2) node[midway,sloped]{\tikz\draw[-{Stealth[length=4pt]}](0,0)--(0.01,0);};
\draw[->] (P2) -- ($(P2)+(0.9,-0.9)$);
\draw[densely dotted] ($(P1)!-0.6!($(P1)+(1.9,-1.1)$)$) -- ($(P1)+(0.8,-0.46)$);
\draw[densely dotted] ($(P2)+(0.65,0.38)$) -- ($(P2)+(-0.8,-0.47)$);
\node[lbl] at ($(P1)+(0.3,-0.45)$) {$r$}; \node[lbl] at ($(P2)+(-0.2,-0.45)$) {$r'$};
\end{tikzpicture}}\ntag{D}{239}{leph201}{19}{}{Class XII Part 2 \textperiodcentered{} Ch 9}}
{\opt{$40^\circ-t^2$}{$40^\circ+t^2$}{$50^\circ-t^2$}{$50^\circ+t^2$}}

\QQ{If AB is incident plane wave front then refracted wave front is ($n_2>n_1$)
\fig{\begin{tikzpicture}[scale=0.7]
\foreach \y in {1,0,-1} \draw[->] (0,\y) -- (1.9,\y);
\draw (0.6,1.1) -- (0.6,-1.1) (1.0,1.1) -- (1.0,-1.1) (1.4,1.1) -- (1.4,-1.1);
\draw (1.9,1) -- (3.0,1) (1.9,-1) -- (3.0,-1) (1.9,0)--(2.9,0);
\draw (3.0,0) ellipse (0.3 and 1.0);
\node[lbl] at (1.0,1.35) {A}; \node[lbl] at (1.0,-1.35) {B};
\node[lbl] at (0.2,0.3) {$n_1$}; \node[lbl] at (2.85,-0.4) {$n_2$}; \node[lbl] at (3.8,0.3) {$n_1$};
\node[lbl] at (3.1,1.3) {Lens};
\end{tikzpicture}}\ntag{A}{258}{leph202}{4}{}{Class XII Part 2 \textperiodcentered{} Ch 10}\nsrc{None found in NEET/AIPMT/AIEEE/JEE papers, 2005--2026 (NCERT concept application)}}
{\optfig{\tikz[scale=0.55]{\draw (0,1) .. controls (0.5,0.4) and (0.5,-0.4) .. (0,-1); \draw (0.5,1) .. controls (1.0,0.4) and (1.0,-0.4) .. (0.5,-1); \node[lbl] at (0,1.25){$A'$}; \node[lbl] at (0.6,1.25){$A''$}; \node[lbl] at (0,-1.25){$B'$}; \node[lbl] at (0.6,-1.25){$B''$};}}
{\tikz[scale=0.55]{\draw (0.4,1) .. controls (-0.1,0.4) and (-0.1,-0.4) .. (0.4,-1); \draw (0.9,1) .. controls (0.4,0.4) and (0.4,-0.4) .. (0.9,-1); \node[lbl] at (0.3,1.25){$A'$}; \node[lbl] at (0.95,1.25){$A''$}; \node[lbl] at (0.3,-1.25){$B'$}; \node[lbl] at (0.95,-1.25){$B''$};}}
{\tikz[scale=0.55]{\draw (0,1) -- (0,-1); \draw (0.6,1) -- (0.6,-1); \node[lbl] at (0,1.25){$A'$}; \node[lbl] at (0.6,1.25){$A''$}; \node[lbl] at (0,-1.25){$B'$}; \node[lbl] at (0.6,-1.25){$B''$};}}
{\tikz[scale=0.55]{\draw (0.3,1) -- (-0.1,-1); \draw (0.8,1) -- (0.4,-1); \node[lbl] at (0.3,1.25){$A'$}; \node[lbl] at (0.85,1.25){$A''$}; \node[lbl] at (-0.15,-1.25){$B'$}; \node[lbl] at (0.45,-1.25){$B''$};}}}

\QQ{The total energy carried by the light wave when it travels from a rarer to a non-reflecting and non-absorbing medium\ntag{D}{262}{leph202}{8}{ \textperiodcentered{} Example 10.1}{Class XII Part 2 \textperiodcentered{} Ch 10}}
{\opt{decreases}{remains same}{increases}{either increases or decreases depending upon angle of incidence}}

\QQ{If the radius of first Bohr orbit is $r$, then the radius of the second Bohr orbit will be\ntag{D}{299}{leph204}{10}{}{Class XII Part 2 \textperiodcentered{} Ch 12}}
{\opt{$2r$}{$8r$}{$4r$}{$2\sqrt2\,r$}}

\QQ{Match the following types of nuclei with examples shown\par\smallskip
\begin{tabular}{@{}ll@{\quad}ll@{}}
\multicolumn{2}{@{}l}{\textit{Column-I}} & \multicolumn{2}{l}{\textit{Column-II}}\\[2pt]
(A) & Isotopes & (i) & ${}_3\mathrm{Li}^7$, ${}_4\mathrm{Be}^7$\\
(B) & Isobars & (ii) & ${}_8\mathrm{O}^{18}$, ${}_9\mathrm{F}^{19}$\\
(C) & Isotones & (iii) & ${}_1\mathrm{H}^1$, ${}_1\mathrm{H}^2$
\end{tabular}\ntag{D}{309}{leph205}{4}{}{Class XII Part 2 \textperiodcentered{} Ch 13}}
{\opt{A--iii, B--i, C--ii}{A--ii, B--iii, C--i}{A--i, B--iii, C--ii}{A--iii, B--ii, C--i}}

\QQ{Which of the following statements is incorrect with reference to `Nuclear force'?\ntag{D}{313}{leph205}{8}{}{Class XII Part 2 \textperiodcentered{} Ch 13}}
{\opt{Potential energy is minimum, if the separation between the nucleons is $0.8\unit{fm}$}{Nuclear force becomes attractive for nucleon distances larger than $0.8\unit{fm}$}{Nuclear force becomes repulsive for nucleon distances less than $0.8\unit{fm}$}{Nuclear force is always attractive}}

\QQ{The range of electrical conductivity ($\sigma$) and resistivity ($\rho$) for metals, among the following, is\ntag{D}{324}{leph206}{2}{}{Class XII Part 2 \textperiodcentered{} Ch 14}}
{\opt{$\rho\to10^{-2}$--$10^{-8}\unit{\Omega m}$; $\sigma\to10^{2}$--$10^{8}\unit{S\,m^{-1}}$}{$\rho\to10^{-5}$--$10^{6}\unit{\Omega m}$; $\sigma\to10^{5}$--$10^{-6}\unit{S\,m^{-1}}$}{$\rho\to10^{11}$--$10^{19}\unit{\Omega m}$; $\sigma\to10^{-11}$--$10^{-19}\unit{S\,m^{-1}}$}{$\rho\to10^{2}$--$10^{8}\unit{\Omega m}$; $\sigma\to10^{-2}$--$10^{-8}\unit{S\,m^{-1}}$}}

\QQ{Which of the following statements is correct for an n-type semiconductor?\ntag{D}{331}{leph206}{9}{}{Class XII Part 2 \textperiodcentered{} Ch 14}}
{\opt{The donor energy level lies just below the bottom of the conduction band}{The donor energy level lies closely above the top of the valence band}{The donor energy level lies at the half way mark of forbidden energy gap}{The donor energy level does not exist}}

\QQ{The circuit shown in the figure contains two ideal diodes $D_1$ and $D_2$. If a cell of emf $3\unit{V}$ and negligible internal resistance is connected as shown, then the current through $70\unit{\Omega}$ resistance (in ampere) is
\fig{\begin{circuitikz}[scale=0.65,transform shape]
\draw (0,3) to[empty diode,l=$D_1$,invert] (2,3) to[R,l=$30\,\Omega$] (4.4,3) -- (4.4,0);
\draw (0,1.6) to[empty diode,l=$D_2$] (2,1.6) to[R,l=$30\,\Omega$] (4.4,1.6);
\draw (0,3) -- (0,0) -- (1.6,0);
\draw[line width=0.8pt] (1.6,-0.45) -- (1.6,0.45); \draw[line width=1.6pt] (1.9,-0.22) -- (1.9,0.22);
\node at (1.75,-0.8) {$3\,\mathrm{V}$};
\draw (1.9,0) -- (2,0) to[R,l=$70\,\Omega$] (4.4,0);
\end{circuitikz}}\ntag{A}{335}{leph206}{13}{}{Class XII Part 2 \textperiodcentered{} Ch 14}\nsrc{AIPMT 2012 (Prelims) -- two ideal diodes with a battery, current supplied; also NEET 2016 (Phase II)}}
{\opt{0}{0.01}{0.02}{0.03}}

\QQ{In determining the refractive index of a glass slab using a travelling microscope, the following readings are tabulated.
\begin{stmts}[(\alph*)]
\item Reading of travelling microscope for ink mark $= 5.123\unit{cm}$
\item Reading of travelling microscope for ink mark through glass slab $= 6.123\unit{cm}$
\item Reading of travelling microscope for chalk dust on glass slab $= 8.123\unit{cm}$
\end{stmts}
From the data, the refractive index of a glass slab is\ntag{A}{229}{leph201}{9}{}{Class XII Part 2 \textperiodcentered{} Ch 9}\nsrc{None found in NEET/AIPMT/AIEEE/JEE papers, 2005--2026 (NCERT concept application)}}
{\opt{1.390}{1.500}{1.601}{1.399}}

\QQ{In an experiment to determine the figure of merit of a galvanometer by half deflection method, a student constructed the following circuit.
\fig{\begin{circuitikz}[scale=0.65,transform shape]
\draw (0,1.2) to[generic,l=$R$] (0,3) to[battery1] (2,3) to[switch,l=$K_1$] (4,3) -- (4,1.2) -- (3.0,1.2) to[rmeter,t=G] (1.0,1.2) -- (0,1.2);
\draw (1.0,1.2) -- (1.0,0) to[generic,l_=$S$] (2.4,0) to[switch,l_=$K_2$] (3.0,0) -- (3.0,1.2);
\fill (1.0,1.2) circle (2pt) (3.0,1.2) circle (2pt);
\end{circuitikz}}
He unplugged a resistance of $5200\unit{\Omega}$ in R. When $K_1$ is closed and $K_2$ is open, the deflection observed in the galvanometer is 26 div. When $K_2$ is also closed and a resistance of $90\unit{\Omega}$ is removed in S, the deflection becomes 13 div. The resistance of galvanometer is nearly\ntag{A}{130}{leph104}{24}{}{Class XII Part 1 \textperiodcentered{} Ch 4}\nsrc{JEE Main 2014 (11 Apr, online) -- half-deflection method for resistance and figure of merit of a galvanometer}}
{\opt{$116.0\unit{\Omega}$}{$45.0\unit{\Omega}$}{$103.0\unit{\Omega}$}{$91.6\unit{\Omega}$}}

\QQ{While determining the coefficient of viscosity of the given liquid, a spherical steel ball sinks by a distance $h = 0.9\unit{m}$. The radius of the ball is $r = \sqrt3\times10^{-3}\unit{m}$. The time taken by the ball to sink in three trials are tabulated as follows.\par\smallskip
{\centering\adjustbox{max width=\linewidth}{\begin{tabular}{|c|c|}\hline
Trial No. & Time taken by the ball to fall by $h$ (s)\\\hline
1 & 2.75\\ 2 & 2.65\\ 3 & 2.70\\\hline
\end{tabular}}\par}\smallskip
The difference between the densities of the steel ball and the liquid is $7000\unit{kg\,m^{-3}}$. If $g = 10.0\unit{m\,s^{-2}}$, then the coefficient of viscosity of the given liquid at room temperature is\ntag{A}{192}{keph202}{13}{}{Class XI Part 2 \textperiodcentered{} Ch 9}\nsrc{None found in NEET/AIPMT/AIEEE/JEE papers, 2005--2026 (NCERT concept application)}}
{\opt{$0.28\unit{Pa\,s}$}{$0.14\unit{Pa\,s}$}{$0.14\times10^{-3}\unit{Pa\,s}$}{$14\unit{Pa\,s}$}}

%%BEGIN-ANALYSIS
\end{multicols}
\par\needspace{16\baselineskip}\vspace{14pt}
\begin{center}{\pgfornament[width=3.2cm,color=gilt]{87}}\par
{\large\color{navy}\scshape KCET 2025 \textperiodcentered{} Answer Key}\end{center}
{\small\setlength{\tabcolsep}{1pt}\renewcommand{\arraystretch}{1.5}
\begin{tabularx}{\linewidth}{@{}*{10}{>{\centering\arraybackslash}X}@{}}\toprule
\textbf{1}\,(4) & \textbf{2}\,(4) & \textbf{3}\,(3) & \textbf{4}\,(4) & \textbf{5}\,(2) & \textbf{6}\,(1) & \textbf{7}\,(1) & \textbf{8}\,(4) & \textbf{9}\,(2) & \textbf{10}\,(1)$^{*}$\\
\textbf{11}\,(3) & \textbf{12}\,(2) & \textbf{13}\,(1) & \textbf{14}\,(4) & \textbf{15}\,(2) & \textbf{16}\,(3) & \textbf{17}\,(4) & \textbf{18}\,(3) & \textbf{19}\,(3) & \textbf{20}\,(3)$^{*}$\\
\textbf{21}\,(1) & \textbf{22}\,(2) & \textbf{23}\,(2) & \textbf{24}\,(1) & \textbf{25}\,(2) & \textbf{26}\,(1) & \textbf{27}\,(3) & \textbf{28}\,(1) & \textbf{29}\,(1) & \textbf{30}\,(4)\\
\textbf{31}\,(4) & \textbf{32}\,(3) & \textbf{33}\,(3) & \textbf{34}\,(4) & \textbf{35}\,(1) & \textbf{36}\,(3) & \textbf{37}\,(3) & \textbf{38}\,(1) & \textbf{39}\,(2) & \textbf{40}\,(2)\\
\textbf{41}\,(2) & \textbf{42}\,(1) & \textbf{43}\,(4) & \textbf{44}\,(2) & \textbf{45}\,(2) & \textbf{46}\,(3) & \textbf{47}\,(1) & \textbf{48}\,(3) & \textbf{49}\,(1) & \textbf{50}\,(2)\\
\textbf{51}\,(2) & \textbf{52}\,(3) & \textbf{53}\,(1) & \textbf{54}\,(4) & \textbf{55}\,(1) & \textbf{56}\,(1) & \textbf{57}\,(4) & \textbf{58}\,(2)$^{*}$ & \textbf{59}\,(4)$^{*}$ & \textbf{60}\,(2)\\\bottomrule\end{tabularx}}\par\smallskip
{\footnotesize Answers solved by the editor; (1)--(4) are the printed options. $^{*}$: see note.\par
\par\smallskip\textbf{Q10.} Intended answer; the given data treat each wheel separately.
\par\smallskip\textbf{Q20.} Field inside a conductor is zero.
\par\smallskip\textbf{Q58.} Lab manual: real depth 3.000\,cm, apparent depth 2.000\,cm.
\par\smallskip\textbf{Q59.} Lab manual: $G = RS/(R-S)$.
}\par
\clearpage
\begin{center}{\pgfornament[width=5cm,color=gilt]{87}}\par
{\fontsize{20}{24}\selectfont\color{navy}\scshape KCET 2025 \textperiodcentered{} Paper Analysis}\par
{\color{gilt}\pgfornament[width=5cm,color=gilt,symmetry=h]{87}}\end{center}
\vspace{4pt}
\begin{multicols}{2}
\textbf{\color{crimson}1. Direct vs Applied}\par\smallskip
{\footnotesize Checked question by question against the four 2026-27 NCERT books (XI Part 1 \& 2, XII Part 1 \& 2). Out of 60, 60 questions are in the present syllabus; percentages are of these.}\par\smallskip
\begin{tabular}{@{}lrr@{}}\toprule
Direct from NCERT & \textbf{38} & 63\%\\
\quad text, summary, table or figure & 34 &\\
\quad NCERT solved Example & 3 &\\
\quad NCERT Exercise & 1 &\\
Applied (concept in NCERT) & \textbf{22} & 37\%\\\midrule
Analysed & 60 &\\
Deleted portion (not analysed) & 0 &\\\midrule
\textbf{Total} & \textbf{60} &\\\bottomrule\end{tabular}\par\smallskip
{\footnotesize Direct: Q 3, 4, 7, 8, 9, 11, 12, 13, 16, 17, 19, 22, 23, 24, 25, 26, 29, 31, 32, 33, 35, 36, 37, 38, 39, 40, 41, 42, 43, 45, 48, 49, 51, 52, 53, 54, 55, 56.\\ From NCERT Examples: Q 9 (Example 6.1), 11 (Example 7.1), 51 (Example 10.1).\\ From NCERT Exercises: Q 33 (Exercise 4.8).}\par\medskip
\textbf{\color{crimson}2. Where the applied questions resemble earlier exams}\par\smallskip
Of the 22 applied questions, \textbf{9} closely resemble a question from an earlier national exam:
NEET/AIPMT/AIIMS 5, JEE Main/Advanced 4, AIEEE 0. The other 13 (Q 1, 2, 5, 6, 10, 20, 28, 34, 44, 46, 50, 58, 60) had no close match in the NEET, AIPMT, AIIMS, AIEEE and JEE papers of 2005--2026 that were searched; they are straightforward applications of NCERT concepts.\par
{\footnotesize Strongest repeats: Q47 (JEE Main 2022, identical wording), Q30 (AIPMT 2015), Q57 (AIPMT 2012), Q59 (JEE Main 2014), Q15 (NEET 2025), Q27 (NEET 2024).}\par\medskip
\textbf{\color{crimson}3. Difficulty level}\par\smallskip
\begin{tabular}{@{}lr@{}}\toprule
Easy & 34\\ Moderate & 26\\ Hard & 0\\\bottomrule\end{tabular}\par\smallskip
{\footnotesize Overall: \textbf{easy to moderate}, slightly easier than 2026; many single-formula questions. The last three questions are practical-based (travelling microscope, half-deflection, viscosity). Class XI: 20 questions, Class XII: 40 (analysed questions only). (Difficulty is the editor's assessment.)}\par\medskip
\columnbreak
\textbf{\color{crimson}4. Chapter-wise distribution}\par\smallskip
{\footnotesize\begin{tabular}{@{}p{4.7cm}c@{\,}r@{\,}r@{\,}r@{}}\toprule
Chapter & Cl. & Qs & Dir & App\\\midrule
Units \& Measurement & XI & 1 & 0 & 1\\
Motion in a Straight Line & XI & 1 & 0 & 1\\
Motion in a Plane & XI & 2 & 2 & 0\\
Laws of Motion & XI & 2 & 0 & 2\\
Work, Energy \& Power & XI & 2 & 2 & 0\\
Rotational Motion & XI & 2 & 1 & 1\\
Gravitation & XI & 2 & 2 & 0\\
Mech.\ Prop.\ of Solids & XI & 1 & 1 & 0\\
Mech.\ Prop.\ of Fluids & XI & 2 & 0 & 2\\
Thermal Prop.\ of Matter & XI & 1 & 0 & 1\\
Thermodynamics & XI & 1 & 1 & 0\\
Kinetic Theory & XI & 1 & 1 & 0\\
Oscillations & XI & 1 & 0 & 1\\
Waves & XI & 1 & 1 & 0\\
Electric Charges \& Fields & XII & 3 & 2 & 1\\
Electrostatic Potential \& Capacitance & XII & 4 & 3 & 1\\
Current Electricity & XII & 5 & 2 & 3\\
Moving Charges \& Magnetism & XII & 6 & 4 & 2\\
Magnetism \& Matter & XII & 1 & 1 & 0\\
Electromagnetic Induction & XII & 2 & 2 & 0\\
Alternating Current & XII & 3 & 2 & 1\\
Electromagnetic Waves & XII & 1 & 1 & 0\\
Ray Optics & XII & 5 & 2 & 3\\
Wave Optics & XII & 2 & 1 & 1\\
Dual Nature & XII & 2 & 2 & 0\\
Atoms & XII & 1 & 1 & 0\\
Nuclei & XII & 2 & 2 & 0\\
Semiconductor Electronics & XII & 3 & 2 & 1\\
\bottomrule\end{tabular}}\par\medskip
\textbf{\color{crimson}5. Teaching takeaways}\par\smallskip
{\footnotesize\begin{itemize}[leftmargin=1em,nosep]
\item NCERT statements asked directly: nuclear-force list (p.\,313), Table 8.1 wavelength ranges (p.\,211), resistivity ranges of metals (p.\,324), donor level just below the conduction band (p.\,331), isotopes/isobars/isotones (p.\,309).
\item Solved Examples repeated: Example 6.1 (centre of mass of three particles on an equilateral triangle), Example 7.1 (perihelion and aphelion), Example 10.1 (energy of light unchanged on refraction).
\item Exercise 4.8 (solenoid with five layers of windings) was asked with new numbers.
\item Three questions come from the Class XII lab manual (travelling microscope, half-deflection method, Stokes' law); revise practical procedures and their calculations.
\end{itemize}}\par\medskip
\textbf{\color{crimson}6. Deleted portions}\par\smallskip
None. Every question of this paper is within the 2026-27 NCERT textbook.\par
\end{multicols}
%%END-ANALYSIS






